
\documentclass{article}

% -----------------------------
% Packages
% -----------------------------
\usepackage{fancyhdr}
\usepackage{extramarks}
\usepackage{amsmath}
\usepackage{amsthm}
\usepackage{amsfonts}
\usepackage{tikz}
\usepackage[plain]{algorithm}
\usepackage{algpseudocode}
\usepackage{amssymb}

\usetikzlibrary{automata,positioning}

% -----------------------------
% Basic Document Settings
% -----------------------------
\topmargin=-0.45in
\evensidemargin=0in
\oddsidemargin=0in
\textwidth=6.5in
\textheight=9.0in
\headsep=0.25in

\linespread{1.1}

\pagestyle{fancy}
\lhead{\hmwkAuthorName}
\chead{\hmwkClass: \hmwkTitle}
\rhead{}
\lfoot{\lastxmark}
\cfoot{\thepage}

\renewcommand\headrulewidth{0.4pt}
\renewcommand\footrulewidth{0.4pt}

\setlength\parindent{0pt}

% -----------------------------
% Problem Section Helpers
% -----------------------------
\newcommand{\enterProblemHeader}[1]{
  \nobreak\extramarks{}{Problem #1 continued on next page\ldots}\nobreak{}
  \nobreak\extramarks{Problem #1 (continued)}{Problem #1 continued on next page\ldots}\nobreak{}
}

\newcommand{\exitProblemHeader}[1]{
  \nobreak\extramarks{Problem #1 (continued)}{Problem #1 continued on next page\ldots}\nobreak{}
  \nobreak\extramarks{Problem #1}{}\nobreak{}
}

\setcounter{secnumdepth}{0}
\newcounter{partCounter}
\newcounter{homeworkProblemCounter}
\setcounter{homeworkProblemCounter}{1}
\nobreak\extramarks{Problem \arabic{homeworkProblemCounter}}{}\nobreak{}

% -----------------------------
% Theorem env (optional)
% -----------------------------
\newtheorem*{theorem}{Theorem}

% -----------------------------
% Homework Problem Environment
%   Usage:
%   \begin{homeworkProblem}            % auto-number
%     % Your answer here
%   \end{homeworkProblem}
%
%   \begin{homeworkProblem}[1.2.3]     % label problem "1.2.3" (e.g. textbook numbering)
%     % Your answer here
%   \end{homeworkProblem}
% -----------------------------
\newenvironment{homeworkProblem}[1][]{
  \def\hwProblemLabel{#1}
  \ifx\hwProblemLabel\empty
    \edef\hwProblemLabel{\arabic{homeworkProblemCounter}}
    \stepcounter{homeworkProblemCounter}
  \fi
  \section{Problem \hwProblemLabel}
  \setcounter{partCounter}{1}
  \enterProblemHeader{\hwProblemLabel}
}{
  \exitProblemHeader{\hwProblemLabel}
}

% -----------------------------
% Homework Details (EDIT THESE)
% -----------------------------
\newcommand{\hmwkTitle}{Homework \#3}
\newcommand{\hmwkClass}{Math 3320}
\newcommand{\hmwkAuthorName}{\textbf{Joseph Quinn}}

% -----------------------------
% Title
% -----------------------------
\title{
  \vspace{2in}
  \textmd{\textbf{\hmwkClass:\ \hmwkTitle}}\\
  \vspace{3in}
}
\author{\hmwkAuthorName}
\date{}

% -----------------------------
% Convenience Macros (optional)
% -----------------------------
\renewcommand{\part}[1]{\textbf{\large Part \Alph{partCounter}}\stepcounter{partCounter}\\}

% Algorithms
\newcommand{\alg}[1]{\textsc{\bfseries \footnotesize #1}}

% Calculus
\newcommand{\deriv}[1]{\frac{\mathrm{d}}{\mathrm{d}x} (#1)}
\newcommand{\pderiv}[2]{\frac{\partial}{\partial #1} (#2)}
\newcommand{\dx}{\mathrm{d}x}

% Statistics
\newcommand{\E}{\mathrm{E}}
\newcommand{\Var}{\mathrm{Var}}
\newcommand{\Cov}{\mathrm{Cov}}
\newcommand{\Bias}{\mathrm{Bias}}

% Proof step (for aligned derivations)
\newcommand{\step}[2]{& #1 & & \text{#2} \\}

% Solution header (optional)
\newcommand{\solution}{\textbf{\large Solution}}

% -----------------------------
% Document
% -----------------------------
\begin{document}

\maketitle
\pagebreak

\begin{homeworkProblem}[A7]
	\begin{enumerate}
	\item[(I1)]
	\begin{align*}
		\max \quad & z = -5x_1 + x_2 + 4x_3     \\
		\text{s.t.} \quad
		           & -3x_1 + 2x_2 - 5x_3 \le 4  \\
		           & x_1 - 4x_2 + x_3 \le -7    \\
		           & -6x_1 - 3x_2 + 8x_3 \le -3 \\
		           & x_1,x_2,x_3 \ge 0
	\end{align*}
	\item[(I2)]
	\begin{align*}
		x_{4} = & 4 +3x_{1}-2x_{2}+5x_{3}  \\
		x_{5} = & -7-x_{1}+4x_{2}-x_{3}    \\
		x_{6} = & -3 +6x_{1}+3x_{2}-8x_{3} \\
		z =     & -5x_{1}+x_{2}+4x_{3}
	\end{align*}

	\item[(I3)]
	Now setting our non basic variables to 0 we get $x_{1}=0,x_{2}=0,x_{3}=0,x_{4}=4,x_{5}=-7,x_{6}=-6$, since we
	have multiple variables that do not satisfy our non negativity contraints this dictionary does not have a basic feasible solution.

	So now lets attempt creating an auxiliary problem to see if we can find starting feasible dictionary.

	\item[(I1)]
	The auxiliary problem will take the following form,
	\begin{align*}
		\max \quad & w = -x_{0}                       \\
		\text{s.t.} \quad
		           & -3x_1 + 2x_2 - 5x_3 -x_{0}\le 4  \\
		           & x_1 - 4x_2 + x_3 -x_{0}\le -7    \\
		           & -6x_1 - 3x_2 + 8x_3 -x_{0}\le -3 \\
		           & x_{0},x_1,x_2,x_3 \ge 0
	\end{align*}
	\item[(I2)]
	This becomes the dictionary,

	\begin{align*}
		x_{4} = & 4 +x_{0}+3x_{1}-2x_{2}+5x_{3}  \\
		x_{5} = & -7+x_{0}-x_{1}+4x_{2}-x_{3}    \\
		x_{6} = & -3+x_{0} +6x_{1}+3x_{2}-8x_{3} \\
		w =    -x_{0}
	\end{align*}

	\item[(I3)]
	There is no basic feasible solution here (yet) so we must move on to the first aux pivot

	\item[(P1)]
	For the aux pivot the first entering variable is always $x_{0}$
	\item[(P2)]
	With $x_{0}$ as our entering variable, the ratio test looks a little different as we need to know what valeus $x_{0}$ can
	take to ensure the constants, and thus basic vairable are non negative.
	\begin{align*}
		x_{0} & \; \text{has no restriction for} \; x_{4} \\
		x_{0} & \ge 7 \; \text{ for } \; x_{5}            \\
		x_{0} & \ge 3 \; \text{ for } \; x_{6}
	\end{align*}

	Meaning $x_{5}$ has the stronger restriction and is our leaving variable.
	\item[(P3)]
	After performing the pivot we get the resulting dictionary,
	\begin{align*}
		x_{4} = & 11 + 4x_{1}-6x_{2}+6x_{3}+x_{5} \\
		x_{0} = & 7 + x_{1} -4x_{2} + x_{3}+x_{5} \\
		x_{6} = & 4 +7x_{1}-x_{2}-7x_{3}+x_{5}    \\
		w =     & -7 -x_{1} +4x_{2}-x_{3}-x_{5}
	\end{align*}

	\item[(P4)]
	Now to get a basic feasible solution we set our non basic variables to 0 resulting in
	$x_{0}=7,x_{1}=0, x_{2}=0, x_{4}=0, x_{4}=11,x_{5}=0,x_{6}=4$ resulting in $z=-7$

	\item[(P1)]
	For the next pivot the only positive coefficient in $w$ is $x_{2}$ so it is our entering variable.
	\item[(P2)]
	Our basic variable restricting are

	\begin{align*}
		x_{2} \le \; \frac{11}{6} \text{ for } x_{4} \\
		x_{2} \le \; \frac{7}{4}\text{ for }  x_{0}  \\
		x_{2} \le \; 4 \text{ for} \; x_{6}          \\
	\end{align*}

	$x_{0}$ has the tightest restriction so it is our leaving variable

	\item[(P3)]
	After performing the pivot we get the resulting dictionary,
	\begin{align*}
		x_{4} = & \frac{1}{2} + \frac{3}{2}x_{0}+\frac{5}{2}x_{1}+\frac{9}{2}x_{3}-\frac{1}{2}x_{5}  \\
		x_{2} = & \frac{7}{4} - \frac{1}{4}x_{0} +\frac{1}{4}x_{1}+\frac{1}{4}x_{3}+\frac{1}{4}x_{5} \\
		x_{6} = & \frac{9}{4}+\frac{1}{4}x_{0}+\frac{27}{4}x_{1}-\frac{29}{4}x_{3}+\frac{3}{4}x_{5}  \\
		w =     & -x_{0}
	\end{align*}

	\item[(P4)]
	Setting our non basic variables to 0 we get $x_{0}=0,x_{1}=0, x_{2}=\frac{7}{4}, x_{3}=0, z_{4}=\frac{1}{2},x_{5}=0, x_{6}=\frac{9}{4}$

	\item[(Aux transition)]
	Now looking at our dictionary we cannot choose an entering variable meaning this is our optimal solution, and since this is an auxiliary problem
	our optimal $w$ value is 0, with $x_{0}=0$ meaning we found a feasible starting dictionary for our original LP problem

	Removing $x_{0}$ and restoring our orignal $z$ equations results in the equivalent dictionary,

	\begin{align*}
		x_{4} = & \frac{1}{2} + \frac{5}{2}x_{1}+\frac{9}{2}x_{3}-\frac{1}{2}x_{5} \\
		x_{2} = & \frac{7}{4} + \frac{1}{4}x_{1}+\frac{1}{4}x_{3}+\frac{1}{4}x_{5} \\
		x_{6} = & \frac{9}{4}+\frac{27}{4}x_{1}-\frac{29}{4}x_{3}+\frac{3}{4}x_{5} \\
		z =     & \frac{7}{4}-\frac{19}{4}x_{1}+\frac{17}{4}x_{3}+\frac{1}{4}x_{5}
	\end{align*}

	Meaning $x_{3}$ is our leaving variable.
	\item[(P1)]
	Looking at our $z$ equation and according to Bland's rule we pick $x_{3}$ to be our entering variable.

	\item[(P2)]
	Via the ratio test we see,

	\begin{align*}
		x_{3} \le \text{ has no restriction for } x_{4}  \\
		x_{3} \le \text{ has no restriction for }  x_{0} \\
		x_{3} \le \; \frac{9}{29}\text{ for} \; x_{6}    \\
	\end{align*}
	\item[(P3)]
	Now after performing the pivot we get the dictionary,

	\begin{align*}
		x_{4} = & \frac{55}{29}+\frac{194}{29}x_{1}-\frac{1}{29}x_{5}-\frac{18}{29}x_{6} \\
		x_{2} = & \frac{53}{29}+\frac{14}{29}x_{1}+\frac{8}{29}x_{5}-\frac{1}{29}x_{6}   \\
		x_{3} = & \frac{9}{29}+\frac{27}{29}x_{1}+\frac{3}{29}x_{5}-\frac{4}{29}x_{6}    \\
		z =     & \frac{89}{29}-\frac{23}{29}x_{1}+\frac{20}{29}x_{5}-\frac{17}{29}x_{6}
	\end{align*}
	\item[(P4)]
	This results in the basic feasible solution $x_{1}=0, x_{2}=\frac{53}{29}, x_{3}=\frac{9}{29}, x_{4}=\frac{55}{29}, x_{5}=0,x_{6}=0$
	with $z= \frac{89}{29}$
	\item[(P1)]
	Using Bland's rule our entering variable is $x_{5}$
	\item [(P2)]
	Via the ratio test we see,

	\begin{align*}
		x_{5} \le \; \frac{55}{194}\text{ for} \; x_{4}  \\
		x_{5} \le \text{ has no restriction for } x_{2}  \\
		x_{5} \le \text{ has no restriction for }  x_{3} \\
	\end{align*}

	Meaning $x_{4}$ is our leaving variable
	\item [(P3)]
	After performing the pivot we get the dictionary,
	\begin{align*}
		x_{5} = & 55 + 194x_{1}-29x_{4}-18x_{6} \\
		x_{2} = & 17 + 54x_{1}-8x_{4}-5x_{6}    \\
		x_{3} = & 6 + 21x_{1}-3x_{4}-2x_{6}     \\
		z =     & 41 + 133x_{1}-20x_{4}-13x_{6}
	\end{align*}
	\item [(P4)]
	this results in the BFS $x_{1}=0,x_{2}=17,x_{3}=4,x_{4}=0,x_{5}=55,x_{6}=0$
	with  $z=41$
	\item[(P1)]
	Via Bland's, $x_{1}$ is our entering variable
	\item[(P2)]
	Via the ratio test we see,

	\begin{align*}
		x_{1} \le \text{ has no restriction for } x_{5}  \\
		x_{1} \le \text{ has no restriction for } x_{2}  \\
		x_{1} \le \text{ has no restriction for }  x_{3} \\
	\end{align*}

	This means there is no restrictions from any equation on our entering variable, meaning we can increase $x_{1}$ arbitrarily large and get an arbitrarily large $z$ value
\\
	\\
	Meaning our  LP is unbounded

\end{homeworkProblem}

\begin{homeworkProblem}[A8]
\end{homeworkProblem}
\end{document}

